\chapter*{Resumen del Proyecto}
\addcontentsline{toc}{chapter}{Resumen del Proyecto}
\markboth{Resumen del Proyecto}{Resumen del Proyecto}

Los mercados financieros no se comportan de manera homogénea a lo largo
del tiempo: alternan periodos de crecimiento sostenido, fases de caída
prolongada y episodios de inestabilidad extrema en los que la volatilidad
se dispara y las correlaciones entre activos se rompen. Estos
``regímenes de mercado'' --- habitualmente descritos de forma informal
como mercados \emph{alcistas} (\textit{bull}), \emph{bajistas}
(\textit{bear}) y de \emph{crisis} --- condicionan de manera directa la
gestión del riesgo, la asignación de activos y la toma de decisiones de
inversión. Identificar de forma objetiva y automática en qué régimen se
encuentra el mercado en cada momento es, por tanto, un problema de gran
interés tanto académico como práctico.

La literatura tradicional ha abordado este problema empleando
principalmente variables técnicas y de precio --- rentabilidades,
volatilidad, medias móviles, momentum --- combinadas con métodos
estadísticos como los modelos de Markov de cambio de régimen
(\textit{regime-switching}) o técnicas de aprendizaje no supervisado como
\textit{K-Means}. Sin embargo, estas variables son, por construcción,
\emph{retrospectivas}: reflejan movimientos de precio que ya se han
producido. En paralelo, el auge de los modelos de lenguaje basados en
\textit{Transformers} y, en particular, de modelos especializados en
lenguaje financiero como \textit{FinBERT}, ha abierto la posibilidad de
extraer de manera sistemática el tono emocional --- positivo, negativo o
neutro --- de grandes volúmenes de noticias financieras, proporcionando
una fuente de información \emph{contemporánea} y potencialmente
complementaria a los indicadores técnicos clásicos.

Este Trabajo de Fin de Grado plantea, diseña e implementa un pipeline
completo de extremo a extremo que combina ambas fuentes de información
--- precio/indicadores técnicos y sentimiento extraído de noticias
financieras --- con el objetivo de identificar regímenes de mercado
mediante técnicas de \textit{clustering} no supervisado, y de evaluar de
forma rigurosa si la incorporación del sentimiento mejora la calidad y la
interpretabilidad económica de los regímenes detectados frente a un
enfoque puramente técnico.

\section*{Objetivos}

El objetivo general del proyecto es \textbf{desarrollar y validar una
metodología de detección de regímenes de mercado basada en técnicas de
\textit{clustering}, comparando un modelo que utiliza únicamente
indicadores técnicos del S\&P 500 frente a modelos que incorporan
información de sentimiento sectorial extraída de noticias financieras
mediante FinBERT}. De este objetivo general se derivan los siguientes
objetivos específicos:

\begin{itemize}
  \item Construir un \textit{dataset} unificado que combine indicadores
  técnicos diarios del índice S\&P 500 con un indicador de sentimiento
  agregado por sector GICS, obtenido a partir del conjunto de noticias
  financieras \textit{FNSPID}.
  \item Diseñar y aplicar un procedimiento reproducible de cálculo de
  sentimiento mediante el modelo \textit{FinBERT}, así como un esquema de
  agregación y suavizado temporal de dicho sentimiento a nivel sectorial.
  \item Construir tres modelos de \textit{clustering} (\textit{K-Means})
  alternativos --- técnico, técnico+sentimiento y sentimiento suavizado
  --- y compararlos mediante métricas internas (coeficiente de
  \textit{silhouette}, análisis de componentes principales) e
  \textit{importancia de variables}.
  \item Incorporar un modelo de \textit{Hidden Markov Model} (HMM) como
  método de comparación secundario, evaluando sus diferencias frente a
  \textit{K-Means} en términos de estabilidad temporal y diferenciación
  económica de los regímenes.
  \item Validar económicamente los regímenes obtenidos mediante el
  cálculo de rentabilidad anualizada, volatilidad, \textit{ratio} de
  Sharpe, \textit{drawdown} y otras métricas, desagregadas por sector
  GICS y por régimen.
\end{itemize}

\section*{Metodología}

El pipeline desarrollado, representado de forma esquemática en la
\autoref{fig:resumen-pipeline}, consta de las siguientes etapas
principales:

\begin{enumerate}
  \item \textbf{Adquisición y filtrado de datos}: se parte del
  \textit{dataset} público \textit{FNSPID} (Financial News and Stock
  Price Integration Dataset), que contiene titulares de noticias
  financieras asociados a tickers del mercado estadounidense. Las
  noticias se filtran para conservar únicamente aquellas asociadas a
  empresas que formaban parte del índice S\&P 500 en la fecha de
  publicación correspondiente, y se reasignan al siguiente día de
  mercado abierto para evitar fugas de información (\textit{look-ahead
  bias}).
  \item \textbf{Cálculo de sentimiento con FinBERT}: cada titular se
  procesa con el modelo \textit{ProsusAI/FinBERT}, obteniéndose las
  probabilidades $p_{\text{pos}}$, $p_{\text{neg}}$ y $p_{\text{neu}}$.
  Se define una puntuación de sentimiento continua
  $s = p_{\text{pos}} - p_{\text{neg}} \in [-1,1]$ y una etiqueta
  categórica (positivo/negativo/neutro) obtenida como el argumento
  máximo de las tres probabilidades.
  \item \textbf{Agregación sectorial}: cada ticker se mapea a su sector
  GICS correspondiente (11 sectores), y para cada combinación
  (fecha, sector) se calculan la media y la desviación típica del
  sentimiento de todas las noticias publicadas. Dado el carácter ruidoso
  de la señal diaria, se aplica adicionalmente una media móvil de 21
  sesiones (\textit{MA21}) que reduce sustancialmente el ruido de alta
  frecuencia mientras preserva las tendencias de medio plazo del
  sentimiento sectorial.
  \item \textbf{Indicadores técnicos}: a partir de los precios diarios
  del S\&P 500 se calculan nueve indicadores técnicos --- RSI, ratios
  respecto a medias móviles de 50 y 200 sesiones, \textit{momentum} a 21
  y 126 sesiones, volatilidad realizada a 30 y 252 sesiones, y
  \textit{skewness} de los retornos a 30 y 252 sesiones.
  \item \textbf{Construcción del \textit{dataset} de clustering}: se
  combinan los indicadores técnicos y las variables de sentimiento
  sectorial (\textit{crudas} y suavizadas) en un único \textit{dataset}
  diario, obteniéndose un total de 1132 observaciones en el rango
  comprendido entre mayo de 2009 y febrero de 2024.
  \item \textbf{\textit{Clustering} con K-Means}: se entrenan tres
  modelos --- \textbf{Modelo A} (solo indicadores técnicos), \textbf{Modelo
  B} (indicadores técnicos + sentimiento sectorial crudo) y
  \textbf{Modelo C} (solo sentimiento sectorial suavizado) --- con
  $K=3$, seleccionado mediante el método del codo y el coeficiente de
  \textit{silhouette}. Adicionalmente se entrena un \textbf{modelo HMM}
  (\textit{Gaussian Hidden Markov Model}, covarianza completa, $K=3$
  estados) sobre los mismos conjuntos de variables como método de
  comparación.
  \item \textbf{Etiquetado e interpretación de regímenes}: los
  \textit{clusters} obtenidos se etiquetan como \emph{Bull},
  \emph{Bear} o \emph{Crisis} mediante el análisis de los centroides
  (momentum, volatilidad y RSI medios de cada \textit{cluster}).
  \item \textbf{Validación económica}: para cada combinación de
  sector GICS y régimen se calculan rentabilidad anualizada,
  volatilidad anualizada, \textit{ratio} de Sharpe ($r_f=2\%$),
  \textit{drawdown} máximo, porcentaje de días positivos y los
  retornos diarios extremos (mejor y peor día).
\end{enumerate}

\begin{figure}[H]
  \centering
  % TODO: insertar figura del pipeline completo (diagrama de bloques)
  \fbox{\parbox{0.85\textwidth}{\centering\vspace{3cm}
  Diagrama de bloques del pipeline completo: FNSPID
  $\rightarrow$ filtrado S\&P 500 $\rightarrow$ FinBERT $\rightarrow$
  agregación sectorial GICS (media/std + MA21) $\rightarrow$ indicadores
  técnicos S\&P 500 $\rightarrow$ \textit{dataset} de clustering
  $\rightarrow$ K-Means (Modelos A/B/C) y HMM $\rightarrow$ validación
  económica por sector y régimen.\vspace{3cm}}}
  \caption{Esquema general del pipeline desarrollado en este proyecto.}
  \label{fig:resumen-pipeline}
\end{figure}

\section*{Resultados principales}

El análisis del coeficiente de \textit{silhouette} en un espacio común de
componentes principales (reteniendo el 80\% de la varianza explicada)
muestra que el \textbf{Modelo A (técnico)} obtiene la mejor cohesión
interna (\textit{silhouette} $\approx 0.193$), seguido del \textbf{Modelo
B (técnico + sentimiento)} ($\approx 0.178$) y, por último, del
\textbf{Modelo C (solo sentimiento suavizado)} ($\approx 0.162$). Esto
indica que las variables técnicas siguen siendo las que mejor separan
estructuralmente los datos, y que el sentimiento por sí solo no es
suficiente para definir \textit{clusters} bien diferenciados.

Sin embargo, la \textbf{validación económica} revela un panorama más
matizado. El Modelo B, que incorpora sentimiento sectorial junto con los
indicadores técnicos, produce regímenes con una diferenciación económica
\emph{más marcada} que el Modelo A: el régimen \emph{Bear} del Modelo B
presenta \textit{ratios} de Sharpe negativos en prácticamente todos los
sectores (p.\ ej., $-0.36$ en \textit{Financial Services} o $-0.26$ en
\textit{Energy}), mientras que el régimen \emph{Crisis} concentra, de
forma contraintuitiva, las rentabilidades anualizadas más elevadas (hasta
$+91.9\%$ en \textit{Real Estate}), reflejando episodios de alta
volatilidad con fuertes recuperaciones (\textit{rebounds}). El sector
\textit{Energy} en régimen \emph{Bear} presenta el \textit{drawdown}
máximo más severo de todo el análisis ($-74.9\%$), así como el mejor y el
peor día individual de toda la muestra ($+16.9\%$ y $-24.7\%$,
respectivamente), lo que evidencia la naturaleza extremadamente volátil
de dicho régimen.

El \textbf{Modelo C}, basado exclusivamente en sentimiento suavizado,
produce centroides prácticamente indistinguibles en las variables
técnicas (p.\ ej., el ratio respecto a la media móvil de 50 sesiones varía
entre 1.0086 y 1.0134 entre regímenes), lo que constituye un
\emph{resultado negativo relevante}: el sentimiento agregado por sí solo,
sin información de precio, no logra capturar la dinámica de los regímenes
de mercado con la misma nitidez que los indicadores técnicos.

La comparación con el modelo \textbf{HMM} muestra que, si bien ambos
métodos identifican regímenes cualitativamente similares, \textit{K-Means}
produce una diferenciación económica \emph{mucho más pronunciada} entre
regímenes (\textit{ratios} de Sharpe que oscilan aproximadamente entre
$-3.1$ y $+3.8$), mientras que el HMM genera estados con \textit{ratios}
de Sharpe mucho más moderados (entre $-0.6$ y $+1.5$ aproximadamente),
consistente con su naturaleza de modelo generativo con transiciones
suaves entre estados.

\section*{Conclusiones}

Los resultados obtenidos confirman que los indicadores técnicos
constituyen la base más sólida para la detección de regímenes de mercado
mediante \textit{clustering}, pero que la incorporación de sentimiento
sectorial extraído mediante FinBERT \textbf{aporta valor incremental en
términos de diferenciación económica} de los regímenes resultantes, en
particular para distinguir episodios de \emph{Bear market} sostenido
frente a episodios de \emph{Crisis} de alta volatilidad con fuertes
oscilaciones. Estos hallazgos abren la puerta a líneas de trabajo futuras
orientadas a la combinación dinámica de ambas fuentes de información, a
la extensión del análisis a otros índices y mercados, y a la exploración
de arquitecturas de \textit{deep learning} secuenciales para la detección
de regímenes.

\vspace{1cm}
\noindent\textbf{Palabras clave:} regímenes de mercado, \textit{clustering},
K-Means, \textit{Hidden Markov Model}, análisis de sentimiento, FinBERT,
noticias financieras, FNSPID, S\&P 500, validación económica.

\clearpage
